\chapter{Diseño de la interfaz de usuario}
\label{anexo2:cap:diseno-interfaz-usuario}
En este capítulo se presenta el diseño de la interfaz de usuario de la
aplicación, que se ha realizado de manera iterativa y centrada en el usuario.

\section{Introducción}
\label{anexo2:sec:diseno-introduccion}
Tanto la interfaz como la experiencia de usuario son aspectos fundamentales en
el desarrollo de aplicaciones y videojuegos. Un diseño adecuado no solo mejora
la usabilidad, sino que también potencia la satisfacción del usuario, lo que es
crucial para el éxito de cualquier producto digital.

Para esta aplicación, el diseño busca ofrecer una experiencia de usuario clara,
motivadora y sencilla, donde la página web sirva como punto de acceso y centro
de gestión y estadísticas, mientras que el videojuego es el medio principal de
interacción, primando la jugabilidad y la claridad en la interfaz.

El público objetivo de la aplicación son personas que buscan una forma divertida
de hacer ejercicio físico desde casa, sin necesidad de equipamiento
especializado y con un enfoque en la mejora de la salud y el bienestar. Para
esto, un desafío importante es eliminar la fricción que pueda existir entre el
usuario y la interacción con la aplicación, facilitando el acceso y la
comprensión de las funcionalidades.

\section{Metodología y principios}
\label{anexo2:sec:diseno-metodologia}
En esta sección se describen las metodologías y principios de diseño empleados
en el desarrollo de la interfaz de usuario, así como las decisiones que han
justificado el enfoque adoptado.

\subsection{Diseño Centrado en el Usuario (DCU)}
El Diseño Centrado en el Usuario (DCU) es una metodología que pone al usuario en
el centro del proceso de diseño. Se basa en un proceso de diseño iterativo que
trata de maximizar la usabilidad y la experiencia del usuario a través de
pruebas y retroalimentación continua \parencites{UCDIDF2025Mar}{UCDW3C2008Apr}%
{UCDUsal}.

Para conseguirlo, se ha desarrollado el sistema intentando siempre asegurar la
usabilidad y la accesibilidad del mismo, teniendo en cuenta las características
y necesidades del usuario final. Esto incluye el seguimiento de buenas prácticas
de diseño, el uso de patrones ampliamente aceptados y la realización de pruebas
de usuario para validar las decisiones de diseño. Este último aspecto se verá
con detalle en la sección \ref{anexo2:sec:diseno-pruebas-usuario}, mientras que
los dos primeros se desarrollan en las siguientes secciones.

\subsection{Principios de Gestalt de la percepción}
Gestalt es una teoría psicológica y de la percepción que se centra en el estudio
de como los humanos perciben patrones y formas en el entorno más que los
elementos individuales que los componen \parencite{Mather2006Foundations}. Estos
principios ayudan a entender cómo los usuarios perciben y procesan la
información visual, lo que es fundamental para el diseño de interfaces
efectivas. A continuación se presentan algunos de los principios utilizados más
relevantes en el diseño de la interfaz.

\subsubsection{Proximidad}
La proximidad se refiere a la tendencia de los usuarios a agrupar elementos que
están cerca unos de otros (ver figura \ref{anexo2:fig:proximidad}). En el diseño de la
interfaz, esto se puede utilizar para comunicar relaciones entre elementos y
organizar la información de manera clara y coherente
\parencite{Soegaard2015Gestalt}.

\diagrama{
    nombre = {Principio de proximidad de Gestalt},
    fuente = {Dominio público, Wikimedia Commons},
    archivo = {anexo2/imgs/diseno/gestalt_proximidad.png},
    etiqueta = {anexo2:fig:proximidad},
}

Por ejemplo, en la página de estadísticas de la aplicación, los datos relacionados
se agrupan en secciones cercanas: los gráficos relacionados con las actividades
realizadas están agrupados en una sección (ver figura
\ref{anexo2:fig:estadisticas-proximidad}). También se puede ver en el menú principal,
donde los elementos relacionados con la información personal del usuario están
separados de los relacionados con la forma física de este por un separador (ver
figura \ref{anexo2:fig:principal-proximidad}).

\diagrama{
    nombre = {Principio de proximidad en la página de estadísticas},
    archivo = {anexo2/imgs/diseno/estadisticas-proximidad.png},
    etiqueta = {anexo2:fig:estadisticas-proximidad},
}

\diagrama{
    nombre = {Principio de proximidad en el menú principal},
    archivo = {anexo2/imgs/diseno/principal-proximidad.png},
    etiqueta = {anexo2:fig:principal-proximidad},
}

\subsubsection{Semejanza}
El principio de semejanza establece que los elementos que comparten
características visuales, como pueden ser el color, la forma o el tamaño,
tienden a ser percibidos como un grupo o una unidad (ver figura
\ref{anexo2:fig:semejanza}). Este principio se utiliza para crear jerarquías
visuales y comunicar la intención de los elementos de la interfaz
\parencite{Soegaard2015Gestalt}.

\diagrama{
    nombre = {Principio de semejanza de Gestalt},
    fuente = {Dominio público, Wikimedia Commons},
    archivo = {anexo2/imgs/diseno/gestalt_semejanza.png},
    etiqueta = {anexo2:fig:semejanza},
    width = 0.5\textwidth,
}

En el caso de la aplicación, se utilizan en los botones de la interfaz. Aquellos
cuya función es similar comparten el color naranja propio de la aplicación (como
pueden ser los botones de iniciar sesión, registrarse o añadir una forma
física), mientras que acciones destructivas como eliminar la cuenta o cerrar
sesión se destacan con un color rojo para diferenciarlos claramente.

También se utiliza durante el videojuego. En este, se utilizan dos colores, el
naranja y el azul, para destacar la izquierda y la derecha del usuario,
respectivamente. Esto luego es utilizado en el minijuego de explotar globos,
donde los globos de cada color deben ser explotados con la parte del cuerpo
correspondiente: los globos naranjas con la mano izquierda y los azules con la
mano derecha (ver figura \ref{anexo2:fig:globos-semejanza}).

\diagrama{
    nombre = {Principio de semejanza en el minijuego de explotar globos},
    archivo = {anexo2/imgs/diseno/globos.png},
    etiqueta = {anexo2:fig:globos-semejanza},
}

% \subsection{Diez heurísticas de Nielsen}

\subsection{Psicología y teoría del color}
La psicología del color es el estudio de como los colores afectan a las
emociones y comportamientos humanos \parencite{ROOHI2019100298}. Por otro lado,
la teoría del color es el estudio de como los colores interactúan entre sí y
cómo pueden combinarse para crear armonías visuales
\parencite{ColorTheoryHandprint}. Ambas disciplinas son fundamentales en el
diseño de interfaces, una elección adecuada de los colores puede mejorar la
usabilidad, la accesibilidad y la experiencia del usuario, así como evocar
emociones y sensaciones específicas.

Las aplicaciones específicas de la psicología del color y la teoría del color en
el diseño de la interfaz de usuario se verán de manera independiente en las
secciones de la aplicación web y del videojuego.


\section{Diseño de la aplicación web}
\label{anexo2:sec:diseno-aplicacion-web}
La aplicación web es el punto de acceso y centro de gestión de la información
para la aplicación. Su diseño se ha realizado siguiendo los principios
anteriores, con un enfoque en la usabilidad, la accesibilidad y la claridad
visual. A continuación se presentan los aspectos más relevantes del diseño de la
aplicación web.

\subsection{Paleta de colores y tipografía}
La paleta de colores utilizada en la aplicación web es una muy sencilla. El
color principal es el naranja\footnotemark, que se utiliza para los botones, gráficos y
elementos destacados de la interfaz.  Este color se ha elegido por su asociación
con la energía, la creatividad y la motivación \parencites{colorpsychology2015}%
{orangeColorPsychology}, lo que es adecuado para una aplicación centrada en el
ejercicio físico y el bienestar.

\footnotetext{La idea de escoger el color naranja para la aplicación surgió al
darme cuenta de que que aún tenía instalada en el móvil la aplicación Strava
(\url{https://www.strava.com/}).}

Además de este color, el blanco y el negro se utilizan abundantemente para
elementos del fondo y textos generales. El color rojo también hace apariencias
ocasionales para botones con acciones destructivas como se mencionó
anteriormente.

En cuanto a la tipografía, se ha utilizado de fuente principal \tech{Roboto}%
{https://fonts.google.com/specimen/Roboto}, creada por Google, ha probado ser
una fuente muy legible en muchas condiciones \parencite{Shu2024Nov}. Es además
una fuente disponible en la mayoría de dispositivos, evitando crear experiencias
dispares entre usuarios.

\subsection{Composición de la página}
La composición de la página juega un papel crucial en la experiencia que los
usuarios tienen en ella. Se presenta un sistema bastante sencillo: tras iniciar
sesión se accede al área personal, desde la que se puede navegar a las páginas
de estadísticas o a la designada para jugar (ver figura
\ref{anexo2:fig:flujo-interfaz}). La navegación entre estas tres páginas se hace
a través de una barra de navegación lateral con varios iconos, cada uno
representando una de las secciones.

\diagrama{
    nombre = {Flujo de navegación de la aplicación web},
    archivo = {anexo2/imgs/diseno/interfaz.pdf},
    etiqueta = {anexo2:fig:flujo-interfaz},
    width = 1\textwidth,
}

Esta estructura es bastante común en aplicaciones web, lo que es ventajoso,
puesto que los usuarios ya están familiarizados con ella. Como indica
\textcite{Nielsen2018Jan}, los usuarios pasan la mayor parte de su tiempo en
otras aplicaciones, por lo que es importante que la interfaz de nuestra
aplicación siga las convenciones y patrones establecidos para que los usuarios
no tengan que aprender a usarla desde cero.

Otra cosa a mencionar es el abundante uso de espacios \enquote{vacíos} en la
interfaz\footnotemark. Estos espacios no contienen información ni elementos
interactivos, pero son importantes para mantener una carga cognitiva baja y
facilitar la navegación y el uso de los componentes útiles de la aplicación.

\footnotetext{Realmente no hay espacios completamente vacíos en la aplicación.
Todos los fondos de la aplicación contienen una \enquote{cascada de texto} que
informa al usuario en qué página se encuentra. Este texto es de un color gris
para que no distraiga al usuario, pero está presente para que no de la sensación
de que la aplicación está incompleta o \enquote{vacía}.}

\subsubsection{Diseño responsivo}
El diseño responsivo es una técnica de diseño web que permite que la interfaz se
adapte a diferentes tamaños de pantalla y dispositivos, garantizando una experiencia
de usuario óptima en todos ellos \parencite{ResponsiveWebDesign2019}. Uno de los objetivos
de este proyecto era crear una aplicación multiplataforma, que pudiese ser utilizada
tanto en ordenadores de sobremesa como en dispositivos móviles\footnotemark. Se puede
ver un ejemplo de esto en la figura \ref{anexo2:fig:interfaz-móvil}, donde se muestra
la página principal desde un teléfono móvil.

\diagrama{
    nombre = {Página principal de la aplicación web en un teléfono móvil},
    archivo = {anexo2/imgs/diseno/movil.png},
    etiqueta = {anexo2:fig:interfaz-móvil},
    width = 0.4\textwidth,
}

\footnotetext{El uso del videojuego en teléfonos móviles es posible, pero no la
forma recomendada de utilizarla. Debido a la distancia a la que se sitúa el
usuario de la pantalla, es recomendable utilizar un ordenador de sobremesa,
portátil o una tablet. La aplicación web está diseñada para ser utilizada en
cualquier dispositivo.}


\section{Diseño del videojuego}
\label{anexo2:sec:diseno-videojuego}
El videojuego es el medio principal de interacción con la aplicación, por lo que
su diseño es fundamental para la experiencia del usuario. A continuación se
presentan los aspectos más relevantes del diseño del videojuego, siguiendo los
principios y metodologías mencionados anteriormente.

\subsection{Minimalismo}
El minimalismo es un principio de diseño que busca simplificar la interfaz al
eliminar elementos innecesarios y centrarse en lo esencial
\parencite{Fessenden2022Feb}.  En el caso de este videojuego es especialmente
importante, ya que el jugador se situará a varios metros de la pantalla, por lo
que es crucial que la interfaz sea lo más clara y sencilla posible.

Este minimalismo y búsqueda de la claridad se refleja en el diseño del sistema
al completo. En todas las pantallas del videojuego se utiliza un fondo negro con
texto blanco y tamaños de letra grandes, lo que facilita la legibilidad desde la
distancia. La réplica del usuario se forma por puntos de colores brillantes y
uniones de color blanco, lo que permite ver la posición en la que se encuentra
el cuerpo de manera rápida y precisa, incluso a grandes distancias.

De la misma manera, los elementos interactivos del videojuego son
intencionalmente sencillos y claros. Los globos que deben ser explotados en su
minijuego son de colores vivos y diferentes según el tipo de acción que se debe
realizar. Por otro lado, los obstáculos que se deben evitar son de color rojo y
parpadean y hacen sonidos al aparecer para llamar la atención del usuario (ver
figura \ref{anexo2:fig:salto}).

\diagrama{
    nombre = {Jugador saltando para esquivar un obstáculo},
    archivo = {anexo2/imgs/diseno/esquivarsalto.png},
    etiqueta = {anexo2:fig:salto}
}

\subsection{Retroalimentación}
Se entiende por retroalimentación (conocido también como \textit{feedback}) a la
información que se da al usuario cuando realiza alguna acción para confirmarle
que esta ha sucedido. Es un aspecto importante de los videojuegos (tanto que se
conoce como \enquote{jugo} o \textit{game juice} en inglés), puesto que provocan
una respuesta emocional muy satisfactoria en el usuario, lo que mejora la
experiencia de juego y la inmersión \parencites{Irajdoost2023Nov}{GameJuiceReddit}.

En el caso de este videojuego, se ha implementado retroalimentación visual y
auditiva.  Por ejemplo, cuando el usuario explota un globo, este se desintegra
en partículas brillantes y se reproduce un sonido satisfactorio. Otro aspecto en
el que se puede ver es en el movimiento del avatar del usuario, que reacciona de
manera fluida, pero inmediata a los movimientos del usuario, lo que proporciona
una sensación de control satisfactoria e intuitiva: es igual de complejo que
mirarse en un espejo.

\section{Pruebas de usuario}
\label{anexo2:sec:diseno-pruebas-usuario}
Para validar las decisiones de diseño y asegurar que la interfaz cumple con las
necesidades y expectativas de los usuarios, se han realizado pruebas de usuario
a lo largo del proceso de desarrollo. Estas pruebas, aunque bastante informales,
se han centrado en observar cómo usuarios reales\footnotemark interactúan con la
interfaz por sí mismos y recoger su retroalimentación sobre la usabilidad y su
experiencia utilizando el sistema.

\footnotetext{Muchas gracias a mis compañeros de piso y amigos, especialmente a
\textit{Encarna} y \textit{Pablo} por criticar el diseño de mi interfaz
innumerables veces.}

Puesto que para cada una de estas pruebas se ha tomado nota de las opiniones de
los usuarios y se está utilizando Git como sistema de control de versiones, se
puede consultar una de estas iteraciones como ejemplo. En la figura
\ref{anexo2:fig:interfaz-vieja} se puede ver el aspecto de la página de datos del
usuario antes de realizar los cambios.

\diagrama{
    nombre = {Página de datos de usuario en la commit ebc3259},
    archivo = {anexo2/imgs/diseno/datosviejo.png},
    etiqueta = {anexo2:fig:interfaz-vieja},
    width = 1\textwidth,
}

De esta página en concreto se recibieron varias críticas, la principal queja
siendo que había mucho espacio vacío y que la información no estaba bien
organizada (en las capturas de pantalla se está usando una magnificación del
200\%). También se criticó que la sección que informa sobre la forma física
actual era difícil de leer y que la barra lateral de navegación era demasiado
gruesa.

Tras recibir estas críticas se realizaron varios cambios en la interfaz, como
centrar el contenido de la página, reducir el tamaño de la barra lateral de
navegación y reorganizar la información para que fuera más fácil de leer.
También se añadió un efecto en el fondo de la página que indica la sección
actual para reducir el espacio vacío y hacer la interfaz más atractiva. El
resultado de estos cambios se puede ver en la figura \ref{anexo2:fig:interfaz-actual}.

\diagrama{
    nombre = {Página de datos de usuario actualmente},
    archivo = {anexo2/imgs/diseno/datosconforma.png},
    etiqueta = {anexo2:fig:interfaz-actual},
    width = 1\textwidth,
}


